\clearpage
\appendixpage
\begin{appendices}

\section{Introduction to Appendices}\label{app:intro}
Appendices contain supplementary material that would disrupt the flow of the main document. The \texttt{appendix} package provides enhanced control over appendix formatting.

\subsection{Basic Appendix Structure}\label{subsec:basic-appendix}
The basic structure for creating appendices is:

\begin{verbatim}
\begin{appendices}
    \section{First Appendix}
    Content of first appendix.

    \section{Second Appendix}
    Content of second appendix.
\end{appendices}
\end{verbatim}

\section{Appendix Content Types}\label{app:content}

\subsection{Data Tables}\label{subsec:appendix-tables}
Appendices are ideal for large tables that would disrupt the main text.

\subsection{Supplementary Figures}\label{subsec:appendix-figures}
Additional figures that support your conclusions.

\subsection{Code Listings}\label{subsec:appendix-code}
Full code implementations that are referenced in the document.

\subsection{Mathematical Proofs}\label{subsec:appendix-proofs}
Detailed proofs or derivations that would interrupt the flow.

\section{Cross-Referencing Appendices}\label{app:references}
Reference appendices from your main text using standard LaTeX references:

\begin{verbatim}
See Appendix~\ref{app:content} for the complete dataset.
The full proof is provided in Appendix~\ref{app:references}.
\end{verbatim}

\end{appendices}